%&XeLaTeX
%
% Copyright 2020  Boyuan Yang <byang@debian.org>
% All Rights Reserved.
%
% Build HOWTO
% -----------
% * xelatex
% * biber
% * xelatex
% * xelatex
%
\documentclass{beamer}
\usetheme{Madrid}
\usecolortheme{beaver}
\usepackage{ctex}
\usepackage{booktabs}
\usepackage{ulem}
\usepackage{pgfpages}
\usepackage{hyperref}
\usepackage{svg}
\hypersetup{
    colorlinks=true,
}
\usepackage{graphicx}
\usepackage{tikz}
\usetikzlibrary{arrows, shapes, chains, graphs}
\usepackage[
    %style=authoryear-icomp,
    backend=biber,
    %sortlocale=zh_CN,
    url=true,
    %maxbibnames=2,
]{biblatex}
\usepackage{filecontents}

\addbibresource{debianrefs.bib}
% set global background; must in preamble
\usebackgroundtemplate%
{%
  \tikz \node[opacity=0] {\includegraphics[width=\paperwidth]{openlogo-nd.pdf}};
}%
\setbeamertemplate{bibliography item}[triangle]
\setbeamertemplate{note page}[plain]
%\setbeameroption{show notes on second screen=right}
\setbeameroption{hide notes}
\title[Intro to Debian and FLOSS development]{Debian及自由软件简介}
\author[hosiet]{Boyuan Yang \\ \href{mailto:byang@debian.org}}
\institute[Debian]{Debian Project}
\date{\today}
\begin{document}
\begin{frame}
\titlepage
\end{frame}
\begin{frame}{内容概览}
\tableofcontents
\end{frame}


%----------------------
% Presentation contents
%----------------------

\section{问题提出}
\subsection{存在问题}
\begin{frame}{Debian软件现状}{为了安全，已经做到了什么？}
作为开源软件项目，Debian提供的软件有以下的特性保证安全：
\begin{itemize}
\item 使用GPG签名验证源代码可信与否（需要上游支持）；
\item 提供编译选项进行安全增强；
\item 提供经过签名的源码包供用户自行编译使用；
\item 对二进制包所在软件仓库\footnote{默认情况下并没有对单独的二进制.deb包签名，这是一个问题。}签名确保安全性；
\item 部分镜像站点启用HTTPS确保传输安全。
\end{itemize}
\pause
\vskip 1em
存在一个盲点：无法确保从\emph{源代码}到所构建\emph{二进制}软件包的对应关系与安全性。
\end{frame}
\begin{frame}{Debian软件现状}{无法保证构建稳定性，为什么不安全？}
构建过程不确定，就无法确保二进制软件包的安全性。

\begin{enumerate}
\item Debian开发者可能为了某种利益\footnote{Debian官方直到2014年才开始接受仅包含源代码的上传
（SourceOnlyUpload）仅上传源代码可以阻止开发者有意准备含有后门或漏洞的二进制软件包。直到
今天，Debian仓库仍然允许二进制包的上传，且在可预见的未来将一直如此。\cite{debianwikisou}}故意上传带后门的二进制deb包；
\item 攻击实例：CVE-2002-0083 for OpenSSH，单比特翻转导致远程提权漏洞；
\item 编译器自身的问题，如XCodeGhost。
\end{enumerate}
\pause

我们希望构建过程是稳定的，是\emph{可重现}的。
\vfill

\begin{quotation}
The motivation behind "reproducible" builds is to allow verification that no flaws have been introduced during the compilation process.\cite{minidebconf20161113}
\end{quotation}
\end{frame}
\subsection{目标描述}
\begin{frame}{工作目标}{期望效果与解决方案}
期望达到的目标：确定的条件下构建的软件包完全相同（“bit-to-bit identical”）。

\iffalse
\begin{figure}[htbp]
\begin{tikzpicture} [%distance=10em,
  every node/.style = {shape=rectangle, draw, align=center}]
  \graph [grow down, branch right=4cm] {
    {源代码, 环境描述 [fill=blue!20], 软件包元信息} -> 可重复的构建过程 [fill=red!20] -> 目标软件包;
  };
\end{tikzpicture}
\end{figure}
\fi

于是，我们可以在多个不同的场合下对软件进行重新构建，给出的二进制包与构建所得的二进制包互相对照是否相同以确保给出的二进制包与源码包的对应关系。
\end{frame}
\section{工作重点}
\subsection{核心问题}
\subsection{问题分类}
\begin{frame}{可重复构建面临的挑战}{不确定性因素的分类}
要排除不确定性因素，包括以下各个方面：

\vspace{1em}
\begin{enumerate}
\item 时间戳；
\item 时区与地区格式；
\item 不确定的文件排序；
\item 不确定的字典、散列值排序；
\item 用户、用户组、\texttt{umask}、环境变量；
\item 构建路径（build path）；\footnote{在Debian 9中容忍同样的构建路径，在sid中则将其纳入考虑范围内。\vspace{1em}}
\item 其它与环境相关的可变情况。
\end{enumerate}
\end{frame}
\subsection{工作思路}
\begin{frame}{Debian中开展可重复构建工作}{工作思路}
\begin{itemize}
\item {在某个Debian环境中连续两次构建一个软件包，改变如下变量：
  \begin{enumerate}
  \item 时间、日期；
  \item 主机名、域名；
  \item 文件系统（disorderfs?）；
  \item 时区、地区locale；
  \item 用户ID、用户组ID；
  \item 内核版本、CPU厂商；
  \item \texttt{/bin/sh}默认shell；
  \item 各类其它环境变量，等等。
  \end{enumerate}
  }
\item {比较两者构建结果是否相同：
  \begin{itemize}
    \item 如相同，记录结果；
    \item 如不同，使用特定工具（diffoscope）检查构建结果，寻找不同点并找到原因所在，记录并提交缺陷报告；
    \item 与上游一起工作，提交补丁，等等。
  \end{itemize}
}
\end{itemize}
\end{frame}
\subsection{辅助工具}

\begin{frame}[t]{工作成果展示}{制定标准}
{\Large SOURCE\_DATE\_EPOCH：\xout{\sout{独立自主}}制定的环境变量标准}

\vspace{1em}
\begin{itemize}
\item 纯数字的UNIX时间戳；
\item 如果环境变量中给定了这个变量，则将所有自己程序的时间戳设为该变量对应的时间。\cite{debian-reproduciblebuilds-sourcedateepoch}
\item 已经被包括\texttt{gcc}在内的各大项目所接受，从根本上解决了时间戳问题。
\end{itemize}
\end{frame}
\begin{frame}[t]{工作成果展示}{与上游的合作}
团队向各种上游发了各式各样的缺陷报告与补丁。略举数例：

\vspace{1em}
\begin{itemize}
\item \texttt{gcc}的\texttt{\_\_DATE\_\_}与\texttt{\_\_TIME\_\_}宏采用\texttt{SOURCE\_DATE\_EPOCH}来防止时间戳问题；
\item \texttt{python} 3.6开始提供有序字典类型；\footnote{\url{https://mail.python.org/pipermail/python-dev/2016-September/146327.html}}
\item \TeX{}Live 2016开始，除了Lua\TeX{}和原始\TeX{}以外均支持\texttt{SOURCE\_DATE\_EPOCH}；\footnote{\url{https://www.preining.info/blog/2016/06/tex-live-2016-released/}}
\item \texttt{doxygen}时间戳问题。\footnote{\url{https://bugs.debian.org/792201}}
\end{itemize}
\end{frame}
\subsection{工作进度}
\begin{frame}{工作进度展示}{Debian 9(Stretch)完成情况}
\only<1>{%
  即将发布的Debian 9可以保证90\%以上的软件包可靠地重复构建。\cite{debian-reproduciblebuilds-homepage}
  \vskip 1em%
  % 三线表！
  \begin{tabular}{ccccc}%
    \toprule%
    套件 & all & reproducible & unrep & FTBFS \\%
    \midrule%
    testing/amd64 & 24132 & 22090 / 91.5\% & 1557 / 6.4\% & 291 / 1.2\% \\
    testing/i386 & 24132 & 21355 / 88.4\% & 2170 / 8.9\% & 374 / 1.5\% \\%
    testing/armhf & 24132 & 21009 / 87.1\% & 1852 / 7.6\% & 695 / 2.8\% \\%
    sid/amd64 & 25399 & 19182 / 75.5\% & 4501 / 17.7\% & 944 / 3.7\% \\%
    sid/i386 & 25399 & 18632 / 73.3\% & 5054 / 19.8\% & 1153 / 4.5\% \\%
    sid/armhf & 25399 & 18591 / 73.2\% & 4823 / 18.9\% & 1215 / 4.7\% \\%
    \bottomrule%
  \end{tabular}%
}%
\only<2>{%
  {\Large Debian Testing (9) x86\_64平台测试情况}%
  \vfill
  
  \includegraphics[width=0.95\paperwidth]{stats_pkg_state_testing_amd64.png}
}%
\only<3>{%
  {\Large Debian Unstable (sid) x86\_64平台测试情况}%
  \vfill
  
  \includegraphics[width=0.95\paperwidth]{stats_pkg_state_unstable_amd64.png}
}%
\end{frame}
\section{其他项目}
\begin{frame}{其他参与可重复性构建的软件项目}
\begin{columns}
\begin{column}{.5\textwidth}
\begin{itemize}
  \item Arch Linux
  \item Bitcoin
  \item coreboot
  \item Debian
  \item F-Droid
  \item Fedora
  \item FreeBSD
\end{itemize}
\end{column}
\begin{column}{.5\textwidth}
\begin{itemize}

  \item Guix
  \item LEDE
  \item NetBSD
  \item OpenWrt
  \item openSUSE
  \item Tails
  \item Tor Browser
\end{itemize}
\end{column}
\end{columns}

\vfill
\begin{center}
……还有更多！
\end{center}
\end{frame}
\section{未来展望}
\begin{frame}{未来展望}{项目带来的额外优势}
使用持续集成方式探测检查软件包的可重复构建性，还可能带来以下额外好处：
\begin{enumerate}
\item 将\texttt{buildd}的一部分构建任务交给持续集成系统，减轻负担；
\item 承担一部分质保任务，探测软件包regression；
\item 检查出软件在边界条件上可能出现的缺陷；
\item 避免不确定性因素对软件质量的影响（互联网质量、时间戳等）。
\end{enumerate}
\vfill
\end{frame}
\section*{参考}
\begin{frame}[t]{参考}{项目资源}
\begin{itemize}
\item 官方网站：\url{https://reproducible-builds.org}
\item Debian测试情况：\url{https://reproducible.debian.net}
\item Debian Wiki主页：\url{https://wiki.debian.org/ReproducibleBuilds}
\item 项目每周报告：\url{https://reproducible.alioth.debian.org/blog/}
\item 项目历史记录：\url{https://wiki.debian.org/ReproducibleBuilds/History}
\end{itemize}
\end{frame}
\begin{frame}{参考}{进一步阅读}
\nocite{*}
\printbibliography
\end{frame}
\begin{frame}{Q\&A}{问答}
\begin{center}
{\Large 问~答~环~节}
\end{center}
\end{frame}
\begin{frame}{版权声明}{Copyright Notice}
本文档以\href{https://creativecommons.org/licenses/by-sa/4.0/legalcode}{知识共享~署名-相同方式共享4.0协议（国际）}发布。
\vspace{1em}

This document is published under \href{https://creativecommons.org/licenses/by-sa/4.0/legalcode}{Creative Commons Attribution-ShareAlike 4.0 International License}.
\end{frame}

%----------------------
% Bibliography
%----------------------
\begin{filecontents*}{debianrefs.bib}
@online{debianwikisou,
  author = {the Debian Project},
  title = {Source Only Upload},
  year = 2016,
  url = {https://wiki.debian.org/SourceOnlyUpload},
  urldate = {2016-11-23}
}
@online{minidebconf20161113,
  author = {Chris Lamb and
            Holger Levsen},
  title = {Reproducible builds status update},
  year = 2016,
  url = {https://people.debian.org/~lamby/2016-11-13-MiniDebConfCambridge/},
  urldate = {2016-11-23},
}
@online{debian-reproduciblebuilds-homepage,
  author = {the Debian Project},
  title = {Overview of various statistics about reproducible builds},
  year = 2016,
  url = {https://tests.reproducible-builds.org/debian/reproducible.html},
  urldate = {2016-11-24},
}
@online{debian-reproduciblebuilds-sourcedateepoch,
  author = {Chris Lamb},
  title = {SOURCE\_DATE\_EPOCH specification},
  year = 2015,
  url = {https://reproducible-builds.org/specs/source-date-epoch/},
}
\end{filecontents*}

%----------------------
% EOF
%----------------------
\end{document}

